\section{Certifying a common instant}\label{sec:instants}

\subsection{Setting}
Resources $1,\dots,n$, $n\ge2$, have independent finite contracts. Resource $i$
is captured $k_i\ge1$ times; every capture has a closed window $[s,e]$ with
$s<e$ on one monotone client clock. The windows of one resource are ordered and
may touch but not overlap: each ends no later than the next one starts. Its
captures are sequential, so its true read instants increase strictly, also
when two windows share an endpoint. In the
\emph{single-run} model of Theorem~\ref{thm:exact} every capture of $i$ returns
the same target observation. The first window of $i$ is $[s_i^-,e_i^-]$ and the
last $[s_i^+,e_i^+]$; for $k_i=1$ they coincide. Each target is \emph{open}: the
contract lets it become current just before the first capture and be
withdrawn just after the last. Without openness the contract already proves
presence outside the captures, which can only help.

The captures \emph{certify a common instant} if every history consistent with
the contracts and the captures has an instant at which every target is
present: $\forall H\,\exists t$. A \emph{fixed} common instant is one instant,
determined by the capture record, at which every target is present in every
consistent history: $\exists t\,\forall H$.

\begin{example}
Let $A$ have a sound field and be read in $[1,2]$ and $[5,6]$, and let $B$ be
read once in $[3,4]$. Every history has a common instant, at $B$'s read; but
$B$ may be present only near its read instant, which can lie anywhere in
$[3,4]$, so no instant is fixed. \code{certify()} returns certified and
\code{fixed_instant()} returns none.
\end{example}

Certification is Cooper and Marzullo's \emph{definitely}
modality~\cite{cooper1991} for the conjunctive predicate ``every target is
present''. With interval timestamps it corresponds to Stoller's Def modality
over the ``definitely occurred before'' order~\cite{stoller2000}; his Inst
modality asks instead for one global state whose local states definitely
overlapped. For conjunctive predicates over observed local states, Garg and
Waldecker decide \emph{definitely} in polynomial time~\cite{garg1996}, and
for two distinct processes Stoller's conditions take the pairwise form of
Theorem~\ref{thm:exact}(i). What differs here is that the local intervals in
which a target is present are never observed: a resource is seen only at its
captures, and Theorem~\ref{thm:criterion} decides from the contract whether
it stayed in between.

\subsection{The exact test}
\begin{theorem}[Exact test]\label{thm:exact}
\textup{(i)} If every field is sound, the captures certify a common instant iff
$e_i^-\le s_j^+$ for all $i\ne j$. \textup{(ii)} If exactly one resource $r$
has a reusable field and the others are sound, they do iff the others satisfy
(i) among themselves and some capture window of $r$ lies in $[L,H]$, where
$L=\max_{i\ne r}e_i^-$ and $H=\min_{i\ne r}s_i^+$. \textup{(iii)} With two or
more reusable fields they never do.
\end{theorem}
\begin{proof}
\emph{(i), sufficiency.} Let $u_i^-$ and $u_i^+$ be the true first and last
read instants of $i$. By soundness and fidelity, $i$ is present throughout
$[u_i^-,u_i^+]$. For $i\ne j$, $u_i^-\le e_i^-\le s_j^+\le u_j^+$, and
$u_i^-\le u_i^+$; hence $t=\max_iu_i^-\le\min_ju_j^+$ lies in every presence
interval. \emph{Necessity.} If $e_i^->s_j^+$ for some $i\ne j$, read $i$ first
at $e_i^-$ and $j$ last at $s_j^+$, other reads anywhere in their windows; let
$i$ become current within $\delta$ before $e_i^-$ and $j$ be withdrawn within
$\delta$ after $s_j^+$, with $2\delta<e_i^--s_j^+$ (openness). No instant has
both.

\emph{(ii).} A reusable resource can leave and return between any two of its
captures, so its presence can be confined to arbitrarily small neighborhoods
of its read instants. If a window of $r$ lies in $[L,H]$, the instant read in
it lies in $[\max u_i^-,\min u_i^+]$ over the sound resources, by (i), and is a
common instant. Conversely, place the sound resources at the extremes as in
(i), so that their certain common interval is exactly $[L,H]$; read $r$ at a
point of each window outside $[L,H]$, which exists because no window lies
inside it, and confine $r$ to small neighborhoods of those points.

\emph{(iii).} Confine two reusable resources to neighborhoods of their read
instants, choose all those instants distinct, which positive windows allow,
and shrink the neighborhoods until they are disjoint.
\end{proof}

The test needs only the two largest $e_i^-$ and the two smallest $s_j^+$, and a
reusable resource's windows are checked in one pass, so it takes time linear
in the number of captures (\code{_violating_pair}). With instantaneous reads
($s=e$) it is the closed-interval intersection that Spaniol et al.\ use for
inducible coherence~\cite{spaniol2009}; what is added is the window model,
completeness in both directions, the lower bounds below and the treatment of
fields that are not sound. For every refused instance, \code{counterexample()}
builds read instants and presence intervals without a common instant, and
\code{verify_counterexample()} checks them independently. Each reusable read
is chosen among $\sum_ik_i+1$ equally spaced rational points of its window, so
an unused instant always remains; an earlier version offered five fixed points
and failed on six reusable resources in one window, which is now a regression
test, with up to twelve reusable reads checked.

\paragraph{Touching windows.} The proofs use the windows of one resource only
through $u_k\le e_k\le s_{k+1}\le u_{k+1}$ for consecutive captures $k,k+1$,
which touching windows satisfy. The refutations place a read at a window
endpoint; where that endpoint is shared with the neighboring window of the same
resource, the neighboring read moves into the interior of its own window, which
has positive length, so the reads stay strictly ordered and the refutation is
unchanged. \code{counterexample()} does this, and the audit compares the test
with brute force on sequential schedules whose windows touch and on random
touching window sets (Table~\ref{tab:audit}). The replay's 1\,s condition, in
which consecutive windows touch, therefore lies within Theorem~\ref{thm:exact}:
a closed bridge $[e^-,s^+]$ may be a single instant, and with sound fields that
instant is common to every history.

\subsection{Schedules and fixed instants}
\begin{corollary}[Schedules]\label{cor:schedule}
\textup{(i)} Certifying $n\ge2$ resources needs at least $2n-1$ captures; the
order $1,\dots,n-1,n,n-1,\dots,1$ attains this, and its singly read resource $n$
needs no sound field. \textup{(ii)} A fixed common instant needs every resource
captured at least twice, hence $2n$ captures, and no reusable field; with sound
fields the fixed instants are exactly $[\max_ie_i^-,\min_is_i^+]$.
\textup{(iii)} With parallel requests, two rounds of $2n$ captures or three
rounds of $2n-1$ suffice.
\end{corollary}
\begin{proof}
\emph{(i)} Two singly captured resources $i,j$ would need $e_i\le s_j$ and
$e_j\le s_i$, so $s_i<e_i\le s_j<e_j\le s_i$, a contradiction; hence at least
$2(n-1)+1$ captures. In the revisit order no first window ends after another
resource's last window starts, and resource $n$ is read once, so
Theorem~\ref{thm:exact}(ii) admits a reusable field there; a singly read
resource of any class is present at its own read, so the order tolerates a
general field at the center as well.

\emph{(ii)} Let $i$ be captured once in $[s_i,e_i]$ and let $t$ be any instant.
Choose $i$'s read instant in the window but different from $t$, and let $i$ be
present only around it (openness): $t$ is not common. A reusable resource is
handled alike however often it is captured. Conversely, with every field sound
and $k_i\ge2$, each $t\in[L,H]$ with $L=\max_ie_i^-$ and $H=\min_is_i^+$
satisfies $u_i^-\le e_i^-\le L\le t\le H\le s_i^+\le u_i^+$ for every $i$, so $t$
is common in every history. For $t<L$, the resource with $e_i^-=L$ can be read
at $e_i^-$ and appear just before it, after $t$; symmetrically for $t>H$. In the
double collect $1,\dots,n,n,\dots,1$ with sequential windows, $L$ is the end of window $n$ and $H$ the start of window $n+1$, so $L\le H$,
with equality when the two windows touch; when every resource is
revisited, Theorem~\ref{thm:exact}(i) is equivalent to $L\le H$, so certifying
and fixing coincide.

\emph{(iii)} Two rounds (all first captures, then all last captures) use $2n$
captures; three rounds ($n-1$ first captures, the single capture, $n-1$ last
captures) use $2n-1$. Two rounds cannot use $2n-1$: the single capture would
have to lie strictly between them.
\end{proof}

Spaniol et al.\ introduced the inverted revisit order, with the last visited
page downloaded once, under instantaneous downloads, and described ETag or
hash comparison as ``the only 100\% reliable method''~\cite{spaniol2009}. The
corollary shows that this order is optimal, and Theorem~\ref{thm:exact} that
it still certifies under request windows whenever the revisited resources have
sound fields, while the singly read resource needs no continuity guarantee at
all.
Part (ii) prices the saved capture: the revisit order places the common
instant inside the center's window, at an instant that varies between
histories, whereas a double collect~\cite{afek1993} of sound fields with closed bridges
($\max e^-\le\min s^+$) names an interval that every history shares. The
preregistered replay protocol used half-open bridges ($\max e^-<\min s^+$),
which reject a resource whose consecutive windows touch although one common
instant suffices; Section~\ref{sec:replay} reports both.

\subsection{Worst cases and hardness}
In the general model the captures of a resource may return different
observations, all with the target body. For one consistent history of a
resource, its \emph{departure pattern} records whether the resource can be
absent before its first capture, whether its path leaves the target in each
gap between consecutive captures, and whether it can leave after the last
capture. A \emph{worst case} is a consistent history whose pattern dominates
the pattern of every consistent history componentwise. Under openness the
first and last entries always depart; \code{departure_pattern()} also handles
contracts that are not open.

\begin{proposition}[Worst cases and hardness]\label{prop:hard}
Let the captures of a resource return any observations of its target. If every
resource has a worst-case history, certification is decidable in polynomial
time; this includes the single-run sound and reusable cases of
Theorem~\ref{thm:exact} and every contract whose observations identify their
states. In general, deciding certification is coNP-complete, even when every
field is sound, contracts are acyclic and windows are arbitrarily short.
\end{proposition}

\paragraph{The sweep.} Suppose every resource has a worst case. Its
\emph{runs} are the maximal blocks of captures joined by gaps that the worst
case does not leave. A run of captures $a,\dots,b$ with $a<b$ is certainly
present on $[e_a,s_b]$; a run of one capture is present only at its read
instant, which may lie anywhere in its window.

\begin{lemma}\label{lem:sweep}
Certification holds iff \textup{(A)} the fixed runs of all resources share a
point, or \textup{(B)} some resource has a one-capture run whose window lies
inside the union of fixed runs of every other resource.
\end{lemma}
\begin{proof}
\emph{Sufficiency.} A gap that the worst case does not leave is left by no
consistent history, since the worst case dominates them all, so in every
history the resource is present on $[u_a,u_b]\supseteq[e_a,s_b]$. Under (A) a
shared point is common to every history; under (B) the one-capture run's read
instant lies in every other resource's fixed runs. \emph{Necessity.} If neither
holds, let every resource follow its worst case, read each run first at $e_a$
and last at $s_b$, read every one-capture run at a point of its window outside
the other resources' common fixed runs and distinct from all other reads, and
let each departing gap be absent except in small neighborhoods of its two
reads. Presence is then each resource's fixed runs plus neighborhoods of
points, and no point lies in all of them: a fixed point fails (A), and a read
point fails (B) or meets no other read. Finitely many compact sets with empty
intersection keep it empty after the neighborhoods shrink.
\end{proof}

\code{worst_case_pattern()} computes, by forward and backward reachability over
the capture layers, whether some consistent history leaves the target in each
gap; the componentwise maximum is the only possible worst case, and one more
layered search decides whether one history realizes it. A two-layer search over
(state, departed) from every state gives reachability with and without a
departure, so a resource with $k$ captures takes $O(|V|(|V|+|E|)+k|V|^2)$ time.
\code{sweep_certainty()} then decides (A) in one sweep over sorted run
endpoints, $O(M\log M)$ for $M$ captures, and (B) with one binary search per
one-capture window and other resource, $O(nM\log M)$. Sound and reusable
single runs have a worst case (no departing gap, respectively all gaps
departing in one history), and so does every contract whose observations
identify their states, because its gaps are then independent.

\paragraph{Membership in coNP.} A certificate of non-certification gives, for
every capture $j$ of resource $i$, a state $q_{ij}$ that emits the recorded
observation. The verifier checks emission, target membership, reachability of
$q_{i1}$ from an initial state and of $q_{i,j+1}$ from $q_{ij}$. Restricting the
consistent histories to those that pass $q_{ij}$ at capture $j$ fixes the
endpoints of every gap, so paths in different gaps no longer constrain each
other: a departing path is chosen in each gap where one exists, every finite
path is permitted and holding times have no lower bound, so these choices
concatenate. The restricted instance therefore has a worst case, and the
verifier runs the sweep on it and accepts iff the sweep finds no certain
common instant. A history without a common instant yields an accepted
certificate, and an accepted certificate yields such a history by
Lemma~\ref{lem:sweep}. The certificate has $O(K\log|V_{\max}|)$ bits for $K$
captures.

\paragraph{Hardness.} Reduce from CNF satisfiability with clauses
$c_1,\dots,c_q$ over variables $v_1,\dots,v_m$. A clock resource $Z$ has states
$f_0\to z_1\to f_1\to\dots\to z_q\to f_q$ with bodies alternating off and on, a
fresh observation per $z_k$ and initial state $f_0$; it is captured at $4k$ and
$4k+2$ with observation $z_k$, so it is certainly present exactly on
$[4k,4k+2]$. Each variable resource $X_v$ has two disconnected branches $T$ and
$F$ with identical observations $r_0,\dots,r_q$, captured at $4k+3$. In branch
$T$, the segment between $r_{k-1}$ and $r_k$ also has a departing path through
an off state if $v\in c_k$; in branch $F$, if $\lnot v\in c_k$. If the formula is
satisfiable, choose branches by a satisfying assignment and in each slot let a
variable of a true literal depart across the whole slot; $Z$ covers every
other instant, so no common instant exists. Conversely, a history without a
common instant fixes one branch per variable, and since $Z$ is present on every
slot, in every slot some variable departs, which its branch allows only if its
literal satisfies the clause. All contracts are acyclic, and every window may
have any width up to $1/2$, down to a point.

Every field in the construction is sound by Theorem~\ref{thm:criterion},
independently of how often it is captured. Each clock observation $z_k$ occurs
at one state, which cannot be left and re-entered. Each variable observation
$r_k$ occurs at most once on each acyclic branch, a path cannot return to it
after leaving, and the branches are disconnected, so no path reaches the equal
observation on the other branch; no path between equal observations leaves
and returns to the target. The variable resources have no worst case: branch
$T$ can leave in the slots of clauses containing $v$, branch $F$ in those
containing $\lnot v$, and when both literals occur the two patterns are
incomparable. Hardness therefore comes from hidden choice behind identical
observations, not from unsound fields. Conjunctive predicates over observed
local states are polynomial~\cite{garg1996}, general predicates are hard to
detect~\cite{chase1998}, and weak conjunctive detection becomes NP-complete
once local events are only partially ordered~\cite{tarafdar1998}; that result
asks whether some consistent global state exists, whereas here every history
consistent with independent contracts must contain a real-time common instant,
and the uncertainty comes from hidden server states. The sufficient condition
of Lemma~\ref{lem:sweep} and the general hardness do not form a dichotomy: a
general field without a worst case can still be certified with more captures.
Two resources with a contract that permits one return, read $A,B,A,B,A$ in
disjoint windows, have a common instant in every history, since $A$ absent at
both $B$ reads would need two departures; the reference semantics certifies
this instance, while the sweep returns unknown.

\subsection{Verification}
\code{src/certification.py} implements the pairwise test, the fixed-instant
test, the counterexample constructor and an independent checker, the
worst-case pattern and sweep, an exponential reference semantics for small
explicit contracts, and the reduction. \code{scripts/check_certification.py}
writes \code{data/derived/certification-audit.json}; Table~\ref{tab:audit} lists
its checks. They check the implementation and the proofs; they do not
replace the proofs.

\begin{table}[t]
\caption{Checks against brute force in the certification audit.}
\label{tab:audit}
\small
\begin{tabularx}{\columnwidth}{@{}Xr@{}}
\toprule
Check & Instances \\
\midrule
Sequential schedules, $n=2,3$, every sound/reusable assignment, against a
search over read instants (\sScheduleCertified{} certified) & \sScheduleChecks{} \\
Random overlapping windows, same oracle & \sWindowChecks{} \\
Sequential schedules whose windows touch, $n=2,3$, same oracle
(\sTouchingCertified{} certified) & \sTouchingChecks{} \\
Random window sets in which a resource's windows may touch, same oracle &
\sTouchingRandom{} \\
Sequential schedules, $n=1,2,3$, against a search for a fixed instant
(\sFixedCertified{} fix one) & \sFixedChecks{} \\
Explicit fresh-token and content-hash machines against the rule & \sMachineChecks{} \\
CNF instances of the reduction, with positive and point windows & \sCnfChecks{} \\
Random explicit contracts, sweep against the reference semantics
(\sSweepDeterministic{} observation-deterministic) & \sSweepChecks{} \\
Random contracts, worst-case pattern against enumeration & \sPatternChecks{} \\
Small contracts, component criterion against all labelings & \sRepairChecks{} \\
Random finite refinements, Proposition~\ref{prop:repair}(ii) & \sRefinementChecks{} \\
Bounded epochs on the S3 deletion model & \sEpochChecks{} \\
\bottomrule
\end{tabularx}
\end{table}

The audit also confirms the minimum sequential schedules that certify $n=1,2,3,4$
resources (\sMinimumCaptures{} captures) and that fix an instant for
$n=1,2,3$ (\sMinimumFixed{}), counterexamples for 2, 6, 9 and 12 reusable reads
in one window, and the one-return $A,B,A,B,A$ instance.
